% Source: https://artofproblemsolving.com/community/c3h1209339_geometry_problem
% Stable identity: aops-topic-1209339-q01
Draw the line through $D$ parallel to $BF$, meeting $AC$ at $G$.

\begin{asy}
usepackage("amsmath");
size(5cm);
defaultpen(fontsize(10pt));
pair A=(1,4),B=(0,0),C=(6,0); pair D=(B+C)/2, E=(A+D)/2; pair F=extension(B,E,A,C); pair G=extension(D,D+(F-B),A,C);
draw(A--B--C--cycle,linewidth(0.9)); draw(A--D); draw(B--F,blue); draw(D--G,red+dashed);
dot(D);dot(E);dot(F);dot(G); label("$A$",A,N);label("$B$",B,SW);label("$C$",C,SE);label("$D$",D,S);label("$E$",E,W);label("$F$",F,NE);label("$G$",G,NE);
label("\mbox{}",(min(currentpicture,user=true)+max(currentpicture,user=true))/2,black); // keeps the text after the figure black when LaTeX typesets the labels inline
\end{asy}

In triangle $BCF$, $D$ is the midpoint of $BC$ and $DG\parallel BF$, so $G$ is the midpoint of $FC$: $FG=GC$.

In triangle $ADG$, $E$ is the midpoint of $AD$ and $EF\parallel DG$ (both lie along or parallel to $BF$), so $F$ is the midpoint of $AG$: $AF=FG$.

Hence $AF=FG=GC$, and $AF=\dfrac{1}{3}AC$.
